\documentclass[11pt]{article}
\usepackage{mathunal-base}
\mathunalfuente{serif}

\renewcommand{\examtitulo}{Primer Examen Parcial de Álgebra Lineal --- Semestre I de 2019 --- A}
\renewcommand{\examfecha}{10 de junio de 2019}

\newcommand{\vf}{\fbox{V}\ \fbox{F}}
\newcommand{\sino}{\fbox{Sí}\ \fbox{No}}
\newcommand{\casillasino}{\fbox{\phantom{XX}}}

\begin{document}

\examheader

\vspace{4pt}
\noindent
\begin{minipage}[t]{0.62\linewidth}
\textbf{Puntaje. Sólo para uso Oficial}\\[4pt]
\begin{tabular}{|c|c|c|c|c|c|}
\hline
1--9 & 10 & 11 & 12 & \textbf{TOTAL} & \textbf{NOTA} \\
(54) & (10) & (14) & (22) & (100) & (5.0) \\ \hline
 & & & & & \\ \hline
\end{tabular}
\end{minipage}

\vspace{6pt}
\noindent\textbf{Instrucciones:} La duración del examen es de 1 hora y 50 minutos. El examen consta de 12 preguntas en dos hojas impresas por ambos lados; verifique que su examen esté completo. En las preguntas con procedimiento justifique sus respuestas en los espacios asignados. No está permitido sacar hojas en blanco ni ningún tipo de apuntes durante el examen. Verifique que su celular esté apagado y que no esté a la mano. No se permite el uso de calculadora.

\vspace{8pt}
\noindent Nombre: \dotfill\ Cédula \dotfill

\vspace{4pt}
\noindent Profesor: \dotfill\ Grupo \dotfill

\vspace{8pt}
\fbox{\begin{minipage}[t][5cm]{\dimexpr\linewidth-2\fboxsep-2\fboxrule}
\textbf{Espacio para completar respuestas o borrador. Marque claramente si hay respuestas en este espacio.}\\
\textbf{No voltee esta hoja hasta ser AUTORIZADO por el encargado de salón.}
\end{minipage}}

\newpage
\noindent\textbf{\large I. Completación} $[54\,\mathrm{pts} = 9\times 6\,\mathrm{pts}]$

\vspace{4pt}
\noindent En las preguntas 1 a 9 escriba su respuesta en el recuadro correspondiente. \textbf{IMPORTANTE:} En esta sección se califica sólo la respuesta (si es correcta o incorrecta); NO se tiene en cuenta el procedimiento o respuestas parcialmente correctas. Revise su respuesta.

\begin{enumerate}
\item{}[6 pts] Sean $u, v$ vectores en $\mathbb{R}^n$ que satisfacen $||u||=2$, $||v||=1$ y el ángulo entre $u$ y $v$ es $\pi/3$ (radianes). Determine $||u-2v||$. Su respuesta debe estar simplificada.

\vspace{4pt}
\fbox{\begin{minipage}[t][2.5cm]{\dimexpr\linewidth-2\fboxsep-2\fboxrule}
$||u-2v|| = $
\end{minipage}}

\item{}[6 pts $=2\times 3$ pts] Determine en cada caso si el sistema homogéneo tiene soluciones no triviales (Escriba SI o NO en el recuadro al lado de cada sistema; 3 puntos por cada una de las dos partes)

\vspace{6pt}
\noindent
\begin{minipage}[t]{0.46\linewidth}
$\begin{aligned}
4x_1+7x_2-3x_3 &=0\\
3x_1+2x_2-7x_3 &=0\\
x_1+5x_2+4x_3 &=0
\end{aligned}$
\end{minipage}%
\hfill\casillasino\hfill
\begin{minipage}[t]{0.36\linewidth}
$\begin{aligned}
y+4z\qquad\quad &=0\\
x+2y+3z+4w &=0\\
x+y-z\qquad &=0\\
2z\qquad\qquad\quad &=0
\end{aligned}$
\end{minipage}%
\hfill\casillasino

\item{}[6 pts] Para el sistema $\begin{aligned}x+y+z&=0\\ y+2z&=0\end{aligned}$ determine un conjunto de vectores $\mathcal{B}$ tal que $gen(\mathcal{B})$ es igual al conjunto de soluciones del sistema.

\vspace{4pt}
\fbox{\begin{minipage}[t][2.5cm]{\dimexpr\linewidth-2\fboxsep-2\fboxrule}
$\mathcal{B} = $
\end{minipage}}

\item{}[6 pts] Sean $u = \begin{bmatrix}2\\x\\-1\\1\end{bmatrix}$, $v = \begin{bmatrix}1\\2\\y\\1\end{bmatrix}$ y $w = \begin{bmatrix}3\\1\\2\\2\end{bmatrix}$ Determine valores de $x$ y $y$ para que $u,v,w$ sean linealmente dependientes.

\vspace{4pt}
\fbox{\begin{minipage}[t][1.8cm]{\dimexpr\linewidth-2\fboxsep-2\fboxrule}
$x = $ \hspace{2cm} $y = $
\end{minipage}}

\item{}[6 pts] Sean $A = \begin{bmatrix}1 & 1\\ -1 & 0\end{bmatrix}$ y $B = \begin{bmatrix}1 & 1\\ 0 & -1\end{bmatrix}$. Simplifique primero y entonces evalúe la expresión algebraica de matrices
\[
A((A+B)^2 - B^2 - AB)A^{-1}
\]
(la primera respuesta debe ser una expresión algebraica de matrices simplificada; la segunda respuesta debe ser una matriz $2\times 2$).

\vspace{4pt}
\fbox{\begin{minipage}[t][3.5cm]{\dimexpr\linewidth-2\fboxsep-2\fboxrule}
Simplificación:

\vspace{6pt}
Evaluación:
\end{minipage}}

\item{}[6 pts] Sea $A = \begin{bmatrix}3 & 4\\ 2 & 1\end{bmatrix}$. Indique si $\mathrm{rango}(AB)=2$ para cada una de las siguientes posibles matrices $B$. (Escriba SI o NO en el recuadro al lado de cada matriz $B$; se descuentan dos puntos por cada respuesta incorrecta)

\vspace{6pt}
\noindent
\begin{minipage}[t]{0.46\linewidth}
$a)\ B=\begin{bmatrix}0&1\\1&0\end{bmatrix}$\hfill\casillasino

\vspace{6pt}
$b)\ B=\begin{bmatrix}1&1\\0&1\end{bmatrix}$\hfill\casillasino
\end{minipage}%
\begin{minipage}[t]{0.5\linewidth}
$c)\ B=\begin{bmatrix}-1&1\\1&-1\end{bmatrix}$\hfill\casillasino

\vspace{6pt}
$d)\ B=\begin{bmatrix}1&-1\\0&-1\end{bmatrix}$\hfill\casillasino
\end{minipage}

\item{}[6 pts] Evalúe $\alpha = \det\begin{bmatrix}a&0&0&b\\0&c&0&d\\e&0&f&0\\0&g&h&0\end{bmatrix}$.

\vspace{4pt}
\fbox{\begin{minipage}[t][2.5cm]{\dimexpr\linewidth-2\fboxsep-2\fboxrule}
$\alpha = $
\end{minipage}}

\item{}[6 pts] Encuentre la matriz inversa de
\[
A=\begin{bmatrix}0&1&1\\1&0&0\\0&1&-1\end{bmatrix}.
\]

\vspace{4pt}
\fbox{\begin{minipage}[t][3cm]{\dimexpr\linewidth-2\fboxsep-2\fboxrule}
$A^{-1} = $
\end{minipage}}

\item{}[6 pts] Sean $A$ y $B$ dos matrices $10\times 10$ que satisfacen $\det(AB)=5$ y $\det(A)=-1$. Marque las afirmaciones verdaderas entre las siguientes. (Escriba V o F en el recuadro al lado de cada afirmación; se descuentan dos puntos por cada respuesta incorrecta)

\vspace{6pt}
\noindent
\begin{minipage}[t]{0.46\linewidth}
$a)$ $A$ es invertible\hfill\casillasino

\vspace{6pt}
$b)$ $\mathrm{rango}(B)<10$\hfill\casillasino
\end{minipage}%
\begin{minipage}[t]{0.5\linewidth}
$c)$ Columnas de $A$ son L.D.\hfill\casillasino

\vspace{6pt}
$d)$ Filas de $B$ son L.I.\hfill\casillasino
\end{minipage}
\end{enumerate}

\vspace{10pt}
\noindent\textbf{\large II. Solución con Procedimiento} $[46\,\mathrm{pts} = 10+14+22]$

\begin{enumerate}
\setcounter{enumi}{9}
\item{}[10 pts] Sean $v_1,v_2,\dots,v_k$ vectores L.I. en $\mathbb{R}^n$ y sea $v_{k+1}$ un vector en $\mathbb{R}^n$ que no pertenece a $gen(v_1,v_2,\dots,v_k)$. Pruebe que los vectores $v_1,v_2,\dots,v_k,v_{k+1}$ son L.I. \textit{Sugerencia:} Suponga que $v_1,v_2,\dots,v_k,v_{k+1}$ son L.D. y deduzca una contradicción. (Escriba claramente su argumento.)

\vspace{5cm}

\item{}[14 pts] Sean
\[
v_1=\begin{bmatrix}1\\2\\2\end{bmatrix},\ v_2=\begin{bmatrix}-1\\0\\-1\end{bmatrix},\ v_3=\begin{bmatrix}2\\6\\5\end{bmatrix},\ w=\begin{bmatrix}x\\y\\z\end{bmatrix}.
\]
\begin{enumerate}
\item{}[10 pts] ¿Cuáles condiciones deben cumplir $x,y,z$ para que $w\in gen(v_1,v_2,v_3)$? En otras palabras, describa restricciones (ecuaciones) en las entradas $x, y$ y $z$ para que $w\in gen(v_1,v_2,v_3)$.

\vspace{4cm}

\item{}[4 pts] ¿Son $v_1,v_2,v_3$ vectores L.I.? Justifique. Si son L.D., dé una combinación lineal no trivial de ellos que sea igual a $\mathbf{0}$.

\vspace{4cm}
\end{enumerate}

\item{}[22 pts] En la figura se muestra una red de flujos.

\vspace{6pt}
\begin{center}
\begin{tikzpicture}[>=Stealth, scale=1]
  \coordinate (A) at (0,0);
  \coordinate (B) at (2.2,1.4);
  \coordinate (C) at (2.2,-1.4);
  \coordinate (D) at (4.4,0);
  \draw[thick,->] (A) -- node[above,sloped]{$f_1$} (B);
  \draw[thick,->] (A) -- node[above,sloped,pos=0.55]{$f_3$} (D);
  \draw[thick,->] (A) -- node[below,sloped]{$f_2$} (C);
  \draw[thick,->] (B) -- node[above,sloped]{$f_4$} (D);
  \draw[thick,->] (C) -- node[below,sloped]{$f_5$} (D);
  \fill (A) circle (2pt) node[below left]{A};
  \fill (B) circle (2pt) node[above right]{B};
  \fill (C) circle (2pt) node[below right]{C};
  \fill (D) circle (2pt) node[above right]{D};
  \draw[->] (-1.4,0) -- (A) node[midway,above]{200};
  \draw[->] (B) ++(-0.9,0.55) -- ++(1.4,0) node[midway,above]{100};
  \draw[->] (D)++(0.3,0) -- ++(1.4,0) node[midway,above]{200};
  \draw[->] (C)++(0.3,0) -- ++(1.4,0) node[midway,below]{100};
\end{tikzpicture}
\end{center}

\begin{enumerate}
\item{}[8 pts] Establezca un sistema de ecuaciones lineales para determinar los diferentes flujos.

\vspace{4cm}

\item{}[10 pts] Resuelva completamente el sistema obtenido usando eliminación gaussiana o de Gauss-Jordan, determinando variables pivote y libres. \textit{Nota:} En esta parte no considere restricciones sobre los flujos.

\vspace{5cm}

\item{}[4 pts] Asumiendo que todos los flujos deben ser en la dirección indicada por las flechas, es decir $f_i\geq 0$, determine los valores mínimo y máximo posibles del flujo $f_4$.

\vspace{4cm}
\end{enumerate}
\end{enumerate}

\examfooter
\end{document}
